\chapter{Programa en la Raspberry Pi}
\label{cap:programa}

\section{Estructura del programa}

El archivo \texttt{main.py} actúa como punto de entrada del sistema. Su función no es ejecutar directamente los modelos, sino preparar el entorno de trabajo: carga los parámetros de \texttt{config.yaml}, crea las señales y colas compartidas, inicializa los GPIO y pone en marcha los distintos workers. Cada worker se ejecuta en un hilo y permanece bloqueado mientras no tenga una tarea pendiente, lo que evita ciclos de consulta continua y reduce el uso innecesario del procesador.

Los principales hilos corresponden a detección, antispoofing, reconocimiento, telemetría, envío de eventos a Supabase y actualización de la galería facial. El control general queda a cargo del Pipeline Manager, que recibe mensajes de todos estos componentes y decide cuándo puede comenzar un nuevo ciclo, cuándo debe abrirse el portón y cómo recuperar un estado seguro ante un error o una apertura inesperada.

\section{Pipeline de procesamiento}

El procesamiento se organiza como una secuencia dirigida por eventos. Solamente se admite una persona dentro del pipeline en cada instante y la detección se habilita únicamente cuando el final de carrera indica que el portón está cerrado. De esta manera se evita procesar varios candidatos al mismo tiempo o iniciar una nueva autorización mientras el acceso se encuentra abierto.

El núcleo del sistema corresponde a la cadena cámara--detección--antispoofing--reconocimiento--decisión. La información avanza de forma ordenada por las siguientes etapas:

\begin{enumerate}
    \item La cámara conectada a la Raspberry Pi captura continuamente imágenes del entorno.
    \item YuNet localiza el rostro y obtiene los puntos de referencia necesarios para recortarlo y alinearlo.
    \item MiniFASNetV2 analiza la muestra para determinar si corresponde a una persona real o a un posible intento de suplantación.
    \item SFace genera un embedding facial que se compara con la galería de personas autorizadas.
    \item La lógica de decisión clasifica el intento como autorizado, desconocido o spoof y comunica el resultado al bloque de control físico.
\end{enumerate}

Cada una de estas etapas se describe a continuación.

\subsection{Detección del rostro}

La cámara Arducam OV5647, conectada mediante la interfaz CSI y administrada con Picamera2, entrega imágenes de $1280 \times 960$ píxeles. El primer worker utiliza YuNet para localizar un rostro y sus puntos de referencia. Cuando aparece un candidato con una confianza superior al umbral configurado, se deshabilita temporalmente la búsqueda de nuevas personas y se guardan dos muestras: un recorte ampliado para la comprobación de vida y la imagen original junto con los datos geométricos del rostro para el reconocimiento.

\subsection{Comprobación de persona}

La primera muestra es procesada por MiniFASNetV2. Este modelo determina si el rostro corresponde a una persona real o a un posible intento de suplantación, por ejemplo mediante una fotografía o una pantalla. Si el puntaje no supera el umbral de antispoofing, la segunda muestra se descarta, se registra un evento de tipo spoof y el sistema vuelve a habilitar la detección. Si la muestra es considerada real, se habilita la etapa de reconocimiento.

\subsection{Reconocimiento y decisión de acceso}

SFace alinea el rostro a partir de los puntos obtenidos por YuNet y genera un embedding. Se selecciona la mejor coincidencia y se comprueba que alcance el umbral configurado. Si no existe una coincidencia válida, la identidad se rechaza y el sistema vuelve al estado de espera.

Cuando la identidad es autorizada, el Pipeline Manager activa la señal luminosa correspondiente y genera un pulso de 500~ms sobre la salida que comanda la central del portón. A continuación espera que el final de carrera detecte primero la apertura y luego el cierre. Recién después de recuperar el estado de portón cerrado vuelve a habilitarse la detección de un nuevo usuario.

\begin{figure}[H]
    \centering
    \begin{tikzpicture}[node distance=1.3cm]
        \node[cajachica] (yunet) {\textbf{Cámara y YuNet}\\ Detección del rostro};
        \node[cajachica, below=of yunet] (fas) {\textbf{MiniFASNetV2}\\ Comprobación de vida};
        \node[cajachica, below=of fas] (sface) {\textbf{SFace}\\ Reconocimiento facial};
        \node[cajachica, below=of sface] (pm) {\textbf{Pipeline Manager}\\ Decisión, GPIO y portón};
        \node[cajachica, right=2.6cm of sface] (cache) {\textbf{Caché local}\\ IDs, nombres y embeddings};
        \node[cajachica, below=of cache] (supa) {\textbf{Supabase}\\ Usuarios y eventos};
        \draw[flecha] (yunet) -- node[etiqueta]{Candidato detectado} (fas);
        \draw[flecha] (fas) -- node[etiqueta]{Rostro validado} (sface);
        \draw[flecha] (sface) -- node[etiqueta]{Resultado del reconocimiento} (pm);
        \draw[dflecha] (sface) -- node[etiqueta]{Consulta de\\ galería facial} (cache);
        \draw[dflecha] (cache) -- node[etiqueta]{Sincronización} (supa);
    \end{tikzpicture}
    \caption{Pipeline principal ejecutado en la Raspberry Pi y relación con la galería facial.}
    \label{fig:pipeline-rpi}
\end{figure}

El diagrama presenta de forma simplificada las etapas principales del procesamiento y los componentes que intervienen en el reconocimiento. La galería facial almacenada localmente permite realizar las comparaciones sin depender de una consulta remota en cada intento de acceso. La sincronización con Supabase se ejecuta de manera independiente, mientras que el Pipeline Manager recibe los resultados de las distintas etapas y administra las decisiones del sistema.

\subsection{Integración de los modelos}

Los modelos se integraron mediante workers y buffers sincronizados. De esta forma cada etapa trabaja únicamente cuando la anterior produjo una muestra válida y el procesamiento más costoso no se ejecuta sobre todos los cuadros de la cámara. Esta organización también permite aislar errores y ajustar cada modelo sin reescribir el programa completo.

\begin{figure}[H]
    \centering
    % Reemplazar el recuadro por la imagen cuando esté disponible:
    % \includegraphics[width=\textwidth]{Imagenes/pipeline_smartgate.png}
    \fbox{\parbox[c][4cm][c]{0.9\textwidth}{\centering Imagen pendiente: \texttt{pipeline\_smartgate.png}}}
    \caption{Pipeline general implementado en SmartGate.}
    \label{fig:pipeline-general-smartgate}
\end{figure}

La figura~\ref{fig:pipeline-general-smartgate} reúne el camino principal de ejecución: captura y detección, validación de vida, reconocimiento, decisión del Pipeline Manager, control del portón, registro de eventos y sincronización con la página web y Supabase. También representa los buffers y colas utilizados para desacoplar los modelos y evitar que una operación de red bloquee la autorización.

Por claridad, el diagrama no muestra tres mecanismos auxiliares que forman parte de la implementación final. En primer lugar, systemd envuelve el proceso completo e inicia SmartGate al encender la Raspberry Pi (sección~\ref{sec:systemd}). En segundo lugar, la sincronización de rostros mantiene una galería local persistente en \texttt{data/galeria\_rostros.npz}; esta copia se carga antes de intentar conectarse a Supabase y permite reconocer usuarios durante una caída de red (sección~\ref{sec:offline}). Finalmente, si el Supabase Worker no puede enviar un evento, la cola se respalda en \texttt{data/eventos\_pendientes.jsonl} y se reenvía al recuperar la conexión (sección~\ref{sec:eventos}). Estos elementos completan el pipeline operativo aunque se hayan omitido de la figura para conservar su legibilidad.

\section{Concurrencia, estados y seguridad del flujo}

La coordinación entre los hilos se realiza con objetos \texttt{Event}, colas seguras y buffers compartidos. Las señales indican, entre otros estados, si el portón está cerrado, si la detección puede trabajar y si existe una muestra disponible para la siguiente etapa. Los modelos permanecen dormidos hasta recibir su evento, por lo que no compiten constantemente por los recursos de la placa.

El gestor emplea estados lógicos como \texttt{IDLE}, \texttt{PROCESSING}, \texttt{WAIT\_GATE\_OPEN}, \texttt{WAIT\_GATE\_CLOSE} y \texttt{CLEANING}. Si el portón se abre de manera inesperada o una etapa informa un error, se impide el ingreso de nuevos candidatos, se espera a que los tres modelos alcancen un punto seguro y se limpian los buffers antes de reanudar el servicio. Esta normalización evita que una imagen anterior quede asociada por error con el siguiente usuario.

El final de carrera se supervisa mediante \texttt{gpiod}. Un flanco ascendente representa que el portón dejó la posición de cerrado y un flanco descendente indica que volvió a cerrarse. Al arrancar, el nivel del pin se lee directamente, ya que todavía podría no haberse producido ningún flanco. Las funciones asociadas a estas interrupciones solamente actualizan el estado y notifican al gestor; las inferencias se mantienen fuera de la atención del GPIO.

\section{Arranque automático mediante systemd}
\label{sec:systemd}

systemd es el gestor de servicios utilizado por Raspberry Pi OS y otras distribuciones Linux. Se ejecuta durante el arranque del sistema operativo y se encarga de iniciar, detener y supervisar procesos definidos como servicios. En SmartGate se utilizó para que el programa comenzara a funcionar automáticamente cada vez que se encendiera la Raspberry Pi, sin necesidad de abrir una terminal ni ejecutar manualmente \texttt{main.py}.

La unidad \texttt{smartgate.service} se configuró como un servicio de tipo \texttt{simple}. En ella se define la carpeta de trabajo del proyecto y se indica que el comando de inicio es el script \texttt{systemd/iniciar\_smartgate.sh}. La opción \texttt{Restart=on-failure} solicita a systemd que vuelva a lanzar el programa si finaliza debido a un error, mientras que \texttt{RestartSec=5} introduce una espera de cinco segundos antes del reintento. Al quedar asociado a \texttt{default.target} y habilitado mediante \texttt{systemctl -{}-user enable smartgate.service}, el servicio se inicia automáticamente con la sesión del usuario.

El script de inicio espera hasta que Raspberry Pi OS haya creado un entorno gráfico Wayland o X11. Esta precaución evita que OpenCV intente utilizar la pantalla antes de que se encuentre disponible. Luego establece las variables \texttt{XDG\_RUNTIME\_DIR}, \texttt{WAYLAND\_DISPLAY} y \texttt{DISPLAY}, cambia a la carpeta del proyecto y utiliza \texttt{exec} para iniciar, sin procesos intermedios innecesarios, el intérprete del entorno virtual:

\begin{center}
\texttt{.venv/bin/python -u app/main.py}
\end{center}

La opción \texttt{-u} desactiva el almacenamiento temporal de la salida estándar, permitiendo observar inmediatamente los mensajes del programa en el registro del servicio. Para comprobar su estado pueden utilizarse \texttt{systemctl -{}-user status smartgate.service} y \texttt{journalctl -{}-user -u smartgate.service}. La figura~\ref{fig:systemd} resume el proceso completo.

\begin{figure}[H]
    \centering
    \begin{tikzpicture}[node distance=1.3cm]
        \node[caja] (on) {\textbf{Encendido}\\ Raspberry Pi OS};
        \node[caja, below=of on] (sd) {\textbf{systemd}\\ \texttt{smartgate.service}};
        \node[caja, below=of sd] (sc) {\textbf{Script de inicio}\\ Espera Wayland/X11};
        \node[caja, below=of sc] (venv) {\textbf{Entorno virtual}\\ \texttt{app/main.py}};
        \draw[flecha] (on) -- (sd);
        \draw[flecha] (sd) -- node[etiqueta]{ExecStart} (sc);
        \draw[flecha] (sc) -- node[etiqueta]{exec} (venv);
    \end{tikzpicture}
    \caption{Secuencia de arranque automático de SmartGate mediante systemd.}
    \label{fig:systemd}
\end{figure}
